\section{Quadratic Forms}\label{sec:quadratic-forms}
In this section we will study about quadratic forms and its equivalent forms which are easier to study than a general quadratic form. We will give bounds on the deviation of number of solutions from average number of solutions. Before all of that let's see what is a quadratic form.
\begin{Definition}{Quadratic Forms}{}
  \index{quadratic form}%
  A quadratic form over $\bbF_q$ is a degree-$2$ \emph{homogeneous} polynomial or the zero polynomial.
\end{Definition}
If $f\in\bbF_q[X_1,\dots,X_n]$ is a quadratic form. Now if $q$ is odd then for any $i,j\in[n]$ we write $a_{i,j}X_iX_j=\frac12a_{i,j}X_iX_j+\frac12a_{ij}X_jX_i$ which leads to the representation $$f(X_1,\dots, X_n)= \sum_{i,j\in[n]}a_{i,j}X_iX_j\quad\text{for all }a_{i,j}\in\bbF_q\text{ with }a_{i,j}=a_{j,i}.$$
From this formulation we can represent $f$ by a matrix denoted by $C_f$ where $(i,j)^{th}$ entry of $C_f$ is $a_{i,j}$. This matrix is called the \emph{coefficient matrix}\index{coefficient matrix}.
\begin{observation}\label{obs:coeff-matrix-symmetric}
  Since $a_{i,j}=a_{j,i}$ we have $C_f^\top=C_f$.
\end{observation}
So if $X$ denotes the column vector of all variables then $f=X^\top C_fX$. With this representation now we can talk about equivalence of two quadratic forms.
\begin{Definition}{Equivalent Quadratic Forms}{}
  \index{quadratic form!equivalent}%
  Let $f,g\in\bbF_q[X_1,\dots, X_n]$ be two quadratic forms over $\bbF_q$. Then $f$ and $g$ are said to be \emph{equivalent} if $f$ can be transformed into $g$ by a non-singular linear transformation of the variables.
\end{Definition}

So let $f,g\in\bbF_q[X_1,\dots, X_n]$ are two equivalent quadratic forms. Then Let $X$ be the variable vector for $f$ and $\tdX$ be the variable vector for $g$. Then $f=X^\top C_fX$ and $g=\tdX^\top C_g\tdX$. Suppose $M$ is the non-singular linear transformation of variables i.e.\ $\tdX=MX$. For odd $q$ we have $$\tdX^\top C_g\tdX=g=(MX)^\top C_f(MX)= X^\top (M^\top C_f M)X,$$so $C_g=M^\top C_fM$.
\begin{observation}\label{obs:coeff-matrix-equivalent}
  If $f,g\in\bbF_q[X_1,\dots, X_n]$ ($q$ odd) are two quadratic forms then $f$ and $g$ are equivalent if and only if there exists an invertible matrix $M\in\bbF_q^{n\times n}$ such that $C_g=M^\top C_fM$ where $M$ is the non-singular linear transformation. 
\end{observation}
Since $\tdX=MX$. So $M$ gives the one-one correspondence between $\mcV(f(X_1,\dots,X_n)=\alpha)$ and $\mcV(g(X_1,\dots,X_n)=\alpha)$ for any $\alpha\in\bbF_q$. 
\begin{observation}\label{obs:solution-correspondence-equivalent}
  If $f,g\in\bbF_q[X_1,\dots, X_n]$ ($q$ odd) are two quadratic forms then $f$ and $g$ are equivalent by the non-singular linear transformation $M\in\bbF_q^{n\times n}$. Then $\gm\mapsto M\gm$ where $\gm\in\bbF_q^n$ gives an one-one correspondence between $\mcV(f(X_1,\dots,X_n)=\alpha)$ and $\mcV(g(X_1,\dots,X_n)=\alpha)$ for any $\alpha\in\bbF_q$. 
\end{observation}
We will use another terminology to make our life easier. If $f\in\bbF_q[X_1,\dots, X_n]$ is a quadratic form then for any $\alpha\in\bbF_q$ we will say $f$ \emph{represents} $\alpha$ if $f(X_1,\dots, X_n)=\alpha$ has a solution.

To calculate and understand the of solutions of quadratic forms we define the integer-valued function $\vth:\bbF_q\to\bbZ$ by $\vth(\alpha)=-1$ for all $\alpha\in\bbF_q^*$ and $\vth(0)=q-1$. Then we have the following properties of $\vth$ function.\index{theta-function@$\theta$-function}
\begin{lemma}{}{lem:theta-convolution}
  For any finite field $\bbF_q$ we have $$\sum_{\alpha\in\bbF_q}\vth(\alpha)=0,$$and for any $\alpha\in\bbF_q$ and integers $1\leq k\leq m$,
  $$\sum_{\substack{\alpha_1,\dots,\alpha_m\in\bbF_q\\ \alpha_1+\cdots+\alpha_m=\alpha}}\vth(\alpha_1)\cdots \vth(\alpha_k)=\begin{cases}
    0 & \text{if $1\leq k<m$}\\
    \vth(\alpha)\cdot q^{m-1} & \text{if $k=m$.}
  \end{cases}$$
\end{lemma}
\begin{proof}
  Since for all $\alpha\in\bbF_q^*$, $\vth(\alpha)=-1$ we have $\sum_{\alpha\in\bbF_q^*}\vth(\alpha)=-(q-1)$. Therefore $\vth(0)+\sum_{\alpha\in\bbF_q^*}\vth(\alpha)=0$. So we have the first result. 

  Now suppose $k<m$. Then we have \begin{align*}
    \sum_{\substack{\alpha_1,\dots,\alpha_m\in\bbF_q\\ \alpha_1+\cdots+\alpha_m=\alpha}}\prod_{j=1}^k\vth(\alpha_j) & = \sum_{\alpha_1,\dots,\alpha_k\in\bbF_q}\lt(\prod_{j=1}^k\vth(\alpha_j)\rt)\sum_{\substack{\alpha_{k+1},\dots,\alpha_m\in\bbF_q\\ \alpha_{k+1}+\cdots+\alpha_m=\alpha-\alpha_1-\cdots-\alpha_k}}1\\ 
    & = q^{m-k-1}\sum_{\alpha_1,\dots,\alpha_k\in\bbF_q}\prod_{j=1}^k\vth(\alpha_j)\\ 
    & = q^{m-k-q}\prod_{j=1}^k\sum_{\alpha_j\in\bbF_q}\vth(\alpha_j) =0
  \end{align*}
  Now assume $k=m$. Here we will use induction. The case $m=1$ holds trivially. Suppose the formula is shown for $m\geq 1$. Then \begin{align*}
    \sum_{\substack{\alpha_1,\dots,\alpha_{m+1}\in\bbF_q\\ \alpha_1+\cdots+\alpha_{m+1}=\alpha}}\prod_{j=1}^{m+1}\vth(\alpha_j) & = \sum_{\substack{\alpha_1,\dots,\alpha_{m+1}\in\bbF_q\\ \alpha_1+\cdots+\alpha_{m+1}=\alpha}}\prod_{j=1}^{m+1}\vth(\alpha_j)+\sum_{\substack{\alpha_1,\dots,\alpha_{m+1}\in\bbF_q\\ \alpha_1+\cdots+\alpha_{m+1}=\alpha}}\prod_{j=1}^m\vth(\alpha_j)\\ 
    & = \sum_{\substack{\alpha_1,\dots,\alpha_{m+1}\in\bbF_q\\ \alpha_1+\cdots+\alpha_{m+1}=\alpha}}\lt[\prod_{j=1}^m\vth(\alpha_j)\rt]\cdot [\vth(\alpha_{m+1})+1]\\ 
    & = \sum_{\alpha_1,\dots,\alpha_{m}\in\bbF_q}\lt[\prod_{j=1}^m\vth(\alpha_j)\rt]\cdot \lt[\vth\lt(\alpha-\sum_{j=1}^m\alpha_j\rt)+1\rt]\\ 
    & = q\sum_{\substack{\alpha_1,\dots,\alpha_{m}\in\bbF_q\\ \alpha_1+\cdots+\alpha_{m}=\alpha}}\prod_{j=1}^m\vth(\alpha_j)& [\text{If $\alpha_1+\cdots+\alpha_m\neq \alpha$ then the term is $0$}]\\ 
    & = q\cdot q^{m-2}\cdot \vth(\alpha) & [\text{By Induction Hypothesis}]\\ 
    & = q^{m-1}\cdot \vth(\alpha)
  \end{align*}
  so now we have the second result too. So we have the lemma. 
\end{proof}
\subsection{Odd Characteristic Quadratic Forms}
\index{quadratic form!odd characteristic}%
We now resume the study of quadratic forms $f\in\bbF_q[X_1,\dots,X_n]$ where $q$ is odd. We will first show that every quadratic form $f\in\bbF_q[X_1,\dots,X_n]$ where $q$ is odd is equivalent to a \emph{diagonal quadratic form}\index{quadratic form!diagonal} i.e. a quadratic form of type $a_1X_q^2+\cdots+a_nX_n^2$ where $a_i\in\bbF_q$ for all $i\in[n]$. 
\begin{lemma}{Reduction by Separating a Variable}{lem:qf-reduction-by-rep}
  Let $q$ be odd and $f\in\bbF_q[X_1,\dots,X_n]$ be a quadratic form with $n\geq 2$. If $f$ represents $\alpha\in\bbF_q^*$ then $f$ is equivalent to $\alpha\cdot X_1^2+g(X_2,\dots, X_n)$ where $g$ is a quadratic form in $n-1$ variables.
\end{lemma}
\begin{proof}
  Since $f$ represents $\alpha$, $\exs\ (\gm_1,\dots,\gm_n)\in\bbF_q^n$ such that $f(\gm_1,\dots, \gm_n)=\alpha$. Since $\alpha\neq 0$ not all $\gm_i$ are zero. So we can find an invertible matrix $M\in\bbF_q^{n\times n}$ such that in the first column of $M$ the entries are $\gm_1,\dots,\gm_n$. Therefore if we apply the non-singular linear transformation by $M$ to $f$, we obtain a quadratic form in $Y_1,\dots, Y_n$ for which the coefficient of $Y_1^2$ is $f(\gm_1,\dots, \gm_n)=\alpha$. Thus $f$ is equivalent to a quadratic form of the type $$\alpha Y_1^2+2b_2Y_1Y_2+\cdots +2bY_1Y_n+h(Y_2,\dots, Y_n)=\alpha(Y_1+b_2\cdot \alpha^{-1}\cdot Y_2+\cdots b_n\cdot \alpha^{-1}\cdot Y_n)^2+g(Y_2,\dots, Y_n)$$ for some $b_2,\dots, b_n\in\bbF_q$ where $g,h\in\bbF_q[X_1,\dots,X_n]$ are two quadratic forms. So now take the non-singular linear transformation \begin{align*}
    Z_1 & = Y_1+b_2\cdot \alpha^{-1}\cdot Y_2+\cdots +b_n\cdot \alpha^{-1}\cdot Y_n&&\\ 
    Z_i & = Y_i\ \forall\ i\in\{2,\dots, n\}&&
  \end{align*}Then we get a quadratic form of the desired type and this is equivalent to $f$. 
\end{proof}

Applying \QFRedlem repeatedly separates variables one by one, eventually reaching a diagonal form. The only subtlety is ensuring at each step that the remaining form represents a non-zero element of $\bbF_q$; this is handled in the proof of \QFDiagthm below.


\begin{Theorem}{Diagonal Form of Quadratic Forms over $\bbF_q$}{thm:qf-diagonal-odd}
  For any quadratic form $f\in\bbF_q[X_1,\dots,X_n]$ where $q$ is odd, there exists $a_1,\dots, a_n\in\bbF_q$ such that $f$ is equivalent to the diagonal quadratic form $a_1 X_1^2+\cdots +a_nX_n^2$.
\end{Theorem}
\begin{proof}
  We will prove this using induction on the number of variables. If $n=1$ then $f(X_1)=a_{1,1}X_1^2$. Then its already diagonal. So suppose $n\geq 2$ and the result holds for quadratic forms in $(n-1)$ variables. 
  
  Let $f\in\bbF_q[X_1,\dots,X_n]$ be a quadratic form in $n$ variables. Now the theorem is trivial if $f$ is a zero polynomial. So suppose $f$ is a non-zero polynomial. If there exists $i\in[n]$, $a_{i,i}\neq 0$ then $f$ represents $a_{i,i}$ since at the point $X_i=1$ and $X_j=0$ for all $j\in[n]\setminus\{i\}$, $f$ evaluates to be $a_{i,i}$. If for all $i\in[n]$, $a_{i,i}=0$ then $\exs\ i,j\in[n]$, $i\neq j$ such that $a_{i,j}=a_{j,i}\neq 0$. Then $f$ represents $2a_{i,j}$ since at the point $X_i=X_j=1$ and $X_t=0$ for all $t\in[n]\setminus\{i,j\}$, $f$ evaluates to be $2a_{i,j}$. 
  
  So $f$ represents some element $a_1\in\bbF_q^*$ and hence by \QFRedlem we have $f$ is equivalent to a quadratic form $a_1 X_1^2+g(X_2,\dots, X_n)$ where $g$ is a quadratic form in $(n-1)$ variables. Now by induction hypothesis $g$ is equivalent to a diagonal quadratic form $a_2X_2^2+\cdots a_nX_n^2$. Therefore $f$ is equivalent to the diagonal quadratic form $a_1 X_1^2+\cdots +a_nX_n^2$. 
\end{proof}

So now we can only talk about the solutions of diagonal quadratic forms since all general quadratic forms in $\bbF_q[X_1,\dots,X_n]$ are equivalent to a diagonal quadratic form. Now let $f\in\bbF_q[X_1,\dots,X_n]$ is equivalent to $a_1X_1^2+\cdots+a_nX_n^2$. Then some $a_i$'s may be zero. Since non-singular linear transformation preserves ranks, equivalent quadratic forms have coefficient matrices with same rank. Therefore non-singular linear transformation also preserves the non-zero ness of determinant of the coefficient matrix. 
\begin{observation}\label{obs:diagonal-rank-det}
  Let $f\in\bbF_q[X_1,\dots,X_n]$ is a quadratic form is equivalent to $a_1X_1^2+\cdots+a_nX_n^2$. Then $\Rank(C_f)=n-|\{i\in[n]\colon a_i=0\}|$. Similarly, $\det(C_f)=0$ if and only if $\prod_{j=1}^na_i=0$. 
\end{observation}
\begin{observation}\label{obs:equivalent-det-value}
  Let $f,g\in\bbF_q[X_1,\dots,X_n]$ are two equivalent quadratic forms by the non-singular linear transformation $M\in\bbF_q^{n\times n}$. Since by \QFEquivCoeffObs we have $C_g=M^\top C_f M$ and therefore $\det(g)=\det(f)\cdot \det(M)^2$. 
\end{observation}
If the coefficient matrix $C_f$ for any quadratic form $f\in\bbF_q[X_1,\dots,X_n]$ has rank $n$ then we call $f$ to be \emph{non-degenerate}\index{quadratic form!non-degenerate}. We also define \emph{determinant of $f$}\index{quadratic form!determinant}, $\det(f)$ to be the determinant of the coefficient matrix $C_f$. 

The two-variable lemma below is the base case for the induction used in \QFCountOddthm. Once we know $Z(a_1X_1^2+a_2X_2^2=\alpha)$ exactly, the convolution property \QFConvProplem allows us to build up the solution count for $n$ variables.

\begin{lemma}{Solutions of a Two-Variable Diagonal Quadratic Form}{lem:qf-two-var-count}
  For odd $q$, let $\alpha\in\bbF_q$, $a_1,a_2\in\bbF_q^*$, and let $\eta\in\Mq$ be the quadratic character of $\bbF_q$. Then $$Z\lt(a_1X_1^2+a_2X_2^2=\alpha\rt)=q+\vth(\alpha)\cdot \eta(-a_1a_2).$$
\end{lemma}
\begin{proof}
  Now we can break $Z(a_1X_1^2+a_2X_2^2=\alpha)$ into sum over all $\alpha_1,\alpha_2\in\bbF_q$ product of $Z(a_1X_1^2=\alpha_1)$ and $Z(a_2X_2^2=\alpha_2$. So we have 
  \begin{align*}
    Z(a_1X_1^2+a_2X_2^2=\alpha) & = \sum_{\substack{\alpha_1,\alpha_2\in\bbF_q\\ \alpha_1+\alpha_2=\alpha}}Z(a_1X_1^2=\alpha_1)\cdot Z(a_2X_2^2=\alpha_2)\\ 
    & = \sum_{\substack{\alpha_1,\alpha_2\in\bbF_q\\ \alpha_1+\alpha_2=\alpha}} (1+\eta(a_1^{-1}\alpha_1))\cdot (1+\eta(a_2^{-1}\alpha_2))\\ 
    & = \sum_{\substack{\alpha_1,\alpha_2\in\bbF_q\\ \alpha_1+\alpha_2=\alpha}} \lt(1+\eta(a_1^{-1}\alpha_1)+\eta(a_2^{-1}\alpha_2)+\eta(a_1^{-1}a_2^{-1}\alpha_1\alpha_2)\rt)\\ 
    & = q+ \eta(a_1)\sum_{\alpha_1\in\bbF_q}\eta(\alpha_1)+\eta(a_2)\sum_{\alpha_2\in\bbF_q}\eta(\alpha_2)+\eta(a_1a_2)\sum_{\substack{\alpha_1,\alpha_2\in\bbF_q\\ \alpha_1+\alpha_2=\alpha}}\eta(\alpha_1\alpha_2)\\ 
    & = q+\eta(a_1a_2)\sum_{\substack{\alpha_1,\alpha_2\in\bbF_q\\ \alpha_1+\alpha_2=\alpha}}\eta(\alpha_1\alpha_2) \by{\MultCharPropsthm[(i)]}\\ 
    & = q+\eta(a_1a_2)\sum_{\gm\in\bbF_q}\eta(\gm(\alpha-\gm))
  \end{align*}
  Now consider the polynomial $g(X)=\alpha X-X^2$. Then by \EtaQuadPolythm using the function $\vth$ we have $\sum_{\gm\in\bbF_q}\eta(g(\gm))=\vth(\alpha)\eta(-1)$. Therefore we have the lemma. 
\end{proof}
\begin{Theorem}{Solution Count for Non-degenerate Diagonal Quadratic Forms}{thm:qf-count-even-n-odd-q}
  \index{quadratic form!solution count}\index{solution count formula}%
  Let $f$ be a non-degenerate quadratic form in $\bbF_q[X_1,\dots,X_n]$ with $n$ even and $q$ odd. Then for any $\alpha\in\bbF_q$,
  \begin{enumerate}[label=(\roman*)]
    \item If $n$ is even $Z\lt(f(X_1,\dots,X_n)=\alpha\rt)=q^{n-1}+\vth(\alpha)\cdot q^{\nicefrac{(n-2)}{2}}\cdot \eta\lt((-1)^{\nicefrac{n}{2}}\cdot \det(f)\rt).$
    \item if $n$ is odd $Z\lt(f(X_1,\dots, X_n)=\alpha\rt)=q^{n-1}+q^{\nicefrac{(n-1)}{2}}\cdot \eta\lt((-1)^{\nicefrac{(n-1)}{2}}\cdot \alpha\cdot \det(f)\rt).$
  \end{enumerate}
\end{Theorem}
\begin{proof}
  By \QFDiagthm $f$ is equivalent to a diagonal quadratic form $g(X_1,\dots, X_n)=a_1X_1^2+\cdots a_nX_n^2$ by the non-singular linear transformation $M\in\bbF_q^{n\times n}$. By \QFEquivDetObs, we have $\eta(\det(f))=\eta(\det(g))$. Hence it suffices to prove the theorem for $g(X_1,\dots, X_n)=\alpha$. Since $f$ is non-degenerate by \QFDiagRankDetObs, $a_i\neq 0$ for all $i\in[n]$. 
  \begin{enumerate}[label=(\roman*)]
    \item Let $n$ is even. So take $m=\nicefrac{n}2$. Then 
  \begin{align*}
    Z(g(X_1,\dots, X_n)=\alpha) & = \sum_{\substack{\alpha_1,\dots,\alpha_m\in\bbF_q\\ \alpha_1+\cdots+\alpha_m=\alpha}}\prod_{j=1}^mZ\lt(a_{2j-1}X_{2j-1}^2+a_{2j}X_{2j}^2=\alpha_j\rt)\\ 
    & = \sum_{\substack{\alpha_1,\dots,\alpha_m\in\bbF_q\\ \alpha_1+\cdots+\alpha_m=\alpha}} \prod_{j=1}^m \lt(q+\vth(\alpha_j)\cdot\eta(-a_{2j-1}a_{2j})\rt)\by{\QFTwoVarlem}\\ 
    & = \sum_{\substack{\alpha_1,\dots,\alpha_m\in\bbF_q\\ \alpha_1+\cdots+\alpha_m=\alpha}} q^m+ \sum_{\substack{\alpha_1,\dots,\alpha_m\in\bbF_q\\ \alpha_1+\cdots+\alpha_m=\alpha}}\prod_{j=1}^m \vth(\alpha_j)\cdot\eta(-a_{2j-1}a_{2j})\\ 
    & = q^{m-1}\cdot q^m+ \eta\lt((-1)^m\det(g)\rt)\sum_{\substack{\alpha_1,\dots,\alpha_m\in\bbF_q\\ \alpha_1+\cdots+\alpha_m=\alpha}}\prod_{j=1}^m\vth(\alpha_j)\\ 
    & = q^{2m-1}+\eta\lt((-1)^m\det(g)\rt)\cdot \vth(\alpha)\cdot q^{m-1}\by{\QFConvProplem}\\ 
    & = q^{n-1}+\eta\lt((-1)^{\nicefrac{n}2}\det(f)\rt)\cdot \vth(\alpha)\cdot q^{\nicefrac{(n-2)}{2}}
  \end{align*}
  \item Now suppose $n$ is odd. We will use induction on $n$. Now this formula is valid for $n=1$. So suppose $n\geq 3$. Then 
  \begin{align*}
    Z(g(X_1,\dots, X_n)=\alpha) & = \sum_{\substack{\alpha_1,\alpha_2\in\bbF_q\\ \alpha_1+\alpha_2=\alpha}}Z(a_1X_1^2=\alpha_1)\cdot Z(a_2X_2^2+\cdots +a_nX_n^2=\alpha_2)\\ 
    & = \sum_{\substack{\alpha_1,\alpha_2\in\bbF_q\\ \alpha_1+\alpha_2=\alpha}} \lt(1+\eta(a_1^{-1}\alpha_1)\rt)\cdot \lt(q^{n-2}+\eta\lt((-1)^{\nicefrac{n-1}2}a_2\cdots a_n\rt)\cdot \vth(\alpha)\cdot q^{\nicefrac{(n-3)}{2}}\rt)\by{first part}\\ 
    & = \sum_{\substack{\alpha_1,\alpha_2\in\bbF_q\\ \alpha_1+\alpha_2=\alpha}}q^{n-1}+q^{n-2}\cdot \eta(a_1)\sum_{\alpha_1\in\bbF_q}\eta(\alpha_1)\\ 
    & \qquad\qquad\qquad\qquad\qquad\qquad+q^{\nicefrac{(q-3)}{2}}\cdot \eta\lt((-1)^{\nicefrac{n-1}2}a_2\cdots a_n\rt)\sum_{\alpha_2\in\bbF_q}\vth(\alpha_2)\\ 
    &\qquad\qquad\qquad\quad\ \; + q^{\nicefrac{(q-3)}{2}}\cdot \eta\lt((-1)^{\nicefrac{n-1}2}a_2\cdots a_n\rt)\sum_{\substack{\alpha_1,\alpha_2\in\bbF_q\\ \alpha_1+\alpha_2=\alpha}}\vth(\alpha_2)\eta(a_1^{-1}\alpha_1)\\ 
    & = q^{n-1}+ q^{\nicefrac{(q-3)}{2}}\cdot \eta\lt((-1)^{\nicefrac{n-1}2}a_1\cdots a_n\rt)\sum_{\substack{\alpha_1,\alpha_2\in\bbF_q\\ \alpha_1+\alpha_2=\alpha}}\vth(\alpha_2)\eta(\alpha_1)\by{\QFConvProplem}\\ 
    & = q^{n-1}+ q^{\nicefrac{(q-3)}{2}}\cdot \eta\lt((-1)^{\nicefrac{n-1}2}\cdot \det(f)\rt)\sum_{\gm\in\bbF_q}\eta(\gm)\cdot \vth(\alpha-\gm)
  \end{align*}So now we only have to calculate $\sum_{\gm\in\bbF_q}\eta(\gm)\cdot \vth(\alpha-\gm)$. Now $$\sum_{\gm\in\bbF_q}\eta(\gm)\cdot \vth(\alpha-\gm)=\sum_{\gm\in\bbF_q}\eta(\gm)\cdot [\vth(\alpha-\gm)+1]=q\cdot \eta(\alpha)$$
  \end{enumerate}
  So we have the results for both cases. 
\end{proof}

The following alternate proof uses Jacobi sums\index{Jacobi sum!connection to quadratic forms} directly and illustrates how the machinery of the previous chapter applies to count solutions of polynomial equations.


\begin{alternate-proof}
  Again it suffices to prove this theorem for the diagonal quadratic form $g(X_1,\dots, X_n)=a_1X_1^2+\cdots a_nX_n^2$ where $a_i\in\bbF_q^*$ for all $i\in[n]$. Here we will work with the extended definition of $\eta$. Let $\lm_0$ is the trivial multiplicative character and $\lm_1=\eta$. Then 
    \begin{align*}
      Z(g(X_1,\dots, X_n)=\alpha) & = \sum_{\substack{\alpha_1,\dots,\alpha_n\in\bbF_q\\ \alpha_1+\cdots+\alpha_n=\alpha}}\prod_{j=1}^nZ\lt(a_jX_j^2=\alpha_j\rt)\\ 
      & = \sum_{\substack{\alpha_1,\dots,\alpha_n\in\bbF_q\\ \alpha_1+\cdots+\alpha_n=\alpha}}\lt(\lm_1(a_j\alpha_j)+\lm_0(a_j\alpha_j)\rt)\\ 
      & = \sum_{\substack{\alpha_1,\dots,\alpha_n\in\bbF_q\\ \alpha_1+\cdots+\alpha_n=\alpha}}\sum_{i_1,\dots,i_n=0}^1\lm_{i_1}(a_1\alpha_1)\cdots \lm_{i_n}(a_n\alpha_n)\\ 
      & = \sum_{i_1,\dots,i_n=0}^1 \lt(\prod_{j=1}^n\lm_{i_j}(a_j)\rt) \sum_{\substack{\alpha_1,\dots,\alpha_n\in\bbF_q\\ \alpha_1+\cdots+\alpha_n=\alpha}}\prod_{j=1}^n\lm_{i_j}(\alpha_j)\\ 
      & = \sum_{i_1,\dots,i_n=0}^1 \lt(\prod_{j=1}^n\lm_{i_j}(a_j)\rt) \cdot J_{\alpha}\lt(\lm_{i_1},\dots, \lm_{i_n}\rt)
    \end{align*}
    Now if not all but some $\lm_{i_j}$'s are trivial then $J_{\alpha}\lt(\lm_{i_1},\dots, \lm_{i_n}\rt)=0$ by \JacobiCompletethm[(ii)]. So this remains the cases when all are $\lm_{i,j}$'s are trivial or non-trivial. Now if $i_j=0$ for all $j\in[n]$ then by \JacobiCompletethm[(i)], $J_{\alpha}\lt(\lm_{i_1},\dots, \lm_{i_n}\rt)=q^{n-1}$. Therefore by \begin{equation}\label{eq:quad-jacobi-etas}
      Z(g(X_1,\dots, X_n)=\alpha) =q^{n-1}+\eta(\det(f))\cdot J_{\alpha}\underbrace{(\eta,\dots, \eta)}_{n-\text{times}}
    \end{equation}If $\alpha\neq 0$ then above expression gives us \begin{equation}\label{eq:quad-jacobi-nonzero}
      Z(g(X_1,\dots, X_n)=\alpha) =q^{n-1}+\eta(\det(f))\cdot \eta^n(\alpha)\cdot J(\eta,\dots, \eta)
    \end{equation}
    \begin{enumerate}[label=(\roman*)]
      \item Now suppose $n$ is even. Let $\alpha\neq 0$. Since $n$ is even $\eta^n$ is trivial character. So from the equation  \eqref{eq:quad-jacobi-nonzero} we have $$Z(g(X_1,\dots, X_n)=\alpha) =q^{n-1}+\eta(\det(f))\cdot J(\eta,\dots, \eta)$$ Let $\chi\in\Xq$ is a non-trivial additive character of $\bbF_q$. Then by \JGrelthm[(ii)] we have $$J(\eta,\dots, \eta) = -\frac1qG(\eta,\chi)^n = -\frac1q\lt(G^2(\eta,\chi)\rt)^{\nicefrac{n}{2}} = -\frac1q(\eta(-1)q)^{\nicefrac{n}2} = -q^{\nicefrac{(n-2)}{2}}\cdot \eta\lt((-1)^{\nicefrac{n}2}\rt)$$Hence if $\alpha\neq 0$ then we have $$Z(g(X_1,\dots, X_n)=\alpha)=q^{n-1}+\eta(\det(f))\cdot \lt(-q^{\nicefrac{(n-2)}{2}}\cdot \eta\lt((-1)^{\nicefrac{n}2}\rt)\rt)=q^{n-1}-q^{\nicefrac{(q-2)}{2}}\eta\lt((-1)^{\nicefrac{q}{2}}\cdot \det(f)\rt)$$
      Now suppose $\alpha= 0$. Then \eqref{eq:quad-jacobi-etas} gives $$Z(g(X_1,\dots, X_n)=\alpha) =q^{n-1}+\eta(\det(f))\cdot J_{0}(\eta,\dots, \eta)$$Since $n$ is even $\eta^n$ is trivial. Therefore by \JacZeroRellem[(ii)] we have $$J_0\underbrace{(\eta,\dots, \eta)}_{n-\text{times}}=\eta(-1)\cdot (q-1)\cdot J\underbrace{(\eta,\dots\eta)}_{(n-1)-\text{times}}$$So by \JGrelthm[(i)] we have $$J\underbrace{(\eta,\dots\eta)}_{(n-1)-\text{times}}=\frac{G(\eta,\chi)^{n-1}}{G(\eta^{n-1},\chi)}=G^{n-2}(\eta,\chi)=\lt(G^2(\eta,\chi)\rt)^{\nicefrac{(n-2)}{2}}= q^{\nicefrac{(n-2)}{2}}\cdot \eta\lt((-1)^{\nicefrac{(n-2)}{2}}\rt)$$Therefore from \eqref{eq:quad-jacobi-etas} we get $$Z(g(X_1,\dots, X_n)=\alpha)=q^{n-1}+(q-1)\cdot q^{\nicefrac{(n-2)}{2}}\cdot \eta\lt((-1)^{\nicefrac{n}{2}}\cdot \det(f)\rt)$$
      So combining the result of $\alpha\neq 0$ and $\alpha=0$ we get the theorem. 
      \item Now suppose $n$ is odd. Then $\eta^n=\eta$ which is non-trivial. If $\alpha=0$ then by \JacZeroRellem[(i)] we have $J_0(\eta,\dots,\eta)=0$. So assume $\alpha\neq 0$. Then in the equation \eqref{eq:quad-jacobi-nonzero} by using \JGrelthm[(i)] we have
      $$J(\eta,\dots,\eta)=\frac{G(\eta,\chi)^{n}}{G(\eta^{n},\chi)}=G^{n-1}(\eta,\chi)=\lt(G^2(\eta,\chi)\rt)^{\nicefrac{(n-1)}{2}}= q^{\nicefrac{(n-1)}{2}}\cdot \eta\lt((-1)^{\nicefrac{(n-1)}{2}}\rt)$$So in both cases $$Z(g(X_1,\dots, X_n)=\alpha) =q^{n-1}+ q^{\nicefrac{(n-1)}{2}}\cdot \eta\lt((-1)^{\nicefrac{(n-1)}{2}}\cdot \alpha\cdot \det(f)\rt)$$This expression works for $\alpha=0$ because then the second term becomes zero. 
    \end{enumerate}
    Hence we have theorem. 
\end{alternate-proof}

\begin{note}
  For any general degree-$2$ polynomial $f\in\bbF_q[X_1,\dots,X_n]$ there exists $g,h\in\bbF_q[X_1,\dots,X_n]$ where $g$ is a quadratic form and $\deg(h)\leq 1$. Then by some non-singular linear transformation we can transform $f$ into an equivalent diagonal quadratic form. Then we obtain an equation of following form:
  $$\tdf(X_1,\dots, X_n)=a_1X_1^2+\cdots a_kX_k^2+b_1X_1+\cdots b_nX_n$$ where $a_i\in\bbF_q^*$ for all $i\in[k]$, $k\leq n$ and $b_j\in\bbF_q$ for all $j\in[n]$. Hence for any $\alpha\in\bbF_q$, the number of solutions of $f(X_1,\dots, X_n)=\alpha$ is same as number of solutions of $\tdf(X_1,\dots, X_n)=\alpha$. 

  If $k<n$ then without loss of generality we can assume $b_n\neq 0$. Then the number of solutions is $q^{n-1}$ as we can arbitrarily substitute elements for $X_1,\dots, X_{n-1}$ and then value of $X_n$ is uniquely determined. 
  
  If $k=n$ then takes the non-singular linear transformation $X_i=Y_i-b_i(2a_i)^{-1}$ for all $i\in[n]$ gives an equivalent diagonal quadratic form $a_1Y_1^2+\cdots +a_nY_n^2=c$ for some $c\in\bbF_q$. Now we can use \QFCountOddthm to find the number of solutions. 
\end{note}

\subsection{Even Characteristic Quadratic Forms}
\index{quadratic form!even characteristic}%
We will now study quadratic forms and its number of solutions over finite fields $\bbF_q$ with $\Char(\bbF_q)=2$ i.e. $q$ is even. Again like in the case of odd characteristic we will first reduce any general quadratic form to canonical forms and then we will try to find the number of solutions of such canonical forms. 

For even characteristic fields we define \emph{non-degenerate} quadratic forms differently. Let $f\in\bbF_q[X_1,\dots,X_n]$ is a quadratic form in $n$-variables. Then $f$ is called non-degenerate if $f$ is not equivalent to a quadratic form in fewer than $n$ variables.
\begin{note}
  This definition is consistent with the odd-characteristic case: there, if $f$ is equivalent to a form in fewer than $n$ variables, \QFDiagthm and \QFDiagRankDetObs show that $\Rank(C_f)<n$, i.e.\ $\det(f)=0$, which is the standard criterion for degeneracy. So in both characteristics one can reduce to the non-degenerate case by passing to the equivalent form in fewer variables.
\end{note}  
\begin{lemma}{Reduction for Quadratic Forms in even characteristic}{lem:qf-reduction-even-char}
  A non-degenerate quadratic form $f\in\bbF_q[X_1,\dots,X_n]$ with $q$ even and $n\geq 3$ is equivalent to $X_1X_2+g(X_3,\dots, X_n)$ where $g$ is a quadratic form in $n-2$ variables.
\end{lemma}
\begin{proof}
  First we will show that $f$ is equivalent to a quadratic form where coefficient of $X_1^2$ is $0$. So let $$f(X_1,\dots, X_n)=\sum_{i,j\in[n]}a_{i,j}X_iX_j$$ where $a_{i,j}=a_{j,i}\in\bbF_q$. If there exists $i\in[n]$ such that $a_{i,i}=0$ then we can permute the variables so that $X_i$ goes to $X1$ and thus $a_{1,1}=0$. So assume $a_{i,i}\neq 0$ for all $i\in[n]$. Now if $a_{i,j}=0$ for all $i\neq j$ then $$f(X_1,\dots, X_n)=a_{1,1}X_1^2+\cdots +a_{n,n}X_n^2=\lt(a_{1,1}^{\nicefrac{q}2}X_1+\cdots +a_{n,n}^{\nicefrac{q}2}X_n\rt)^2$$then $f$ is equivalent to $Y_1^2$ by the non-singular linear transformation\begin{align*}
    Y_1&=a_{1,1}^{\nicefrac{q}2}X_1+\cdots +a_{n,n}^{\nicefrac{q}2}X_n&&\\
    Y_i&=X_i\ \forall\ i\in\{2,\dots, n\}&&
  \end{align*}This contradicts the non-degenerate property of $f$ \ctr So there exists $i,j\in[n]$, $i\neq j$ such that $a_{i,j}\neq 0$. Without loss of generality suppose $a_{2,3}\neq 0$. Now $$f(X_1,\dots, X_n)=a_{2,2}X_2+X_2(a_{1,2}X_1+a_{2,3}X_3+\cdots +a_{2,n}X_n)+g_1(X_1,X_3\dots, X_n)$$Then take the non-singular linear transformation
  \begin{align*}
    Y_3&=a_{2,3}^{-1}(a_{1,2}X_1+a_{2,3}X_3+\cdots +a_{2,n}X_n)&&\\
    Y_i&=X_i\ \forall\ i\in[n]\setminus\{3\}&&
  \end{align*}This gives an equivalent quadratic form $a_{22}Y_2^2+Y_2Y_3+g_2(Y_1,Y_3,\dots, Y_n)$. Let $b$ is the coefficient of $X_1^2$ in $g_2$. So now take the non-singular linear transformation
  \begin{align*}
    Z_2&=\lt(a_{2,2}^{-1}b\rt)^{\nicefrac{q}2}Y_1+Y_2&&\\
    Z_i&=Y_i\ \forall\ i\in[n]\setminus\{2\}&&
  \end{align*}This transformation makes the coefficient of $Z_1^2$ to be $0$. Now we'll use $f$ to denote this final quadratic form with variables $X_1,\dots, X_n$. Since $f$ is non-degenerate not all $a_{i,j}$'s are $0$ where $i\neq j$. Without loss of generality suppose $a_{1,2}\neq 0$. Then take the non-singular linear transformation
    \begin{align*}
    Y_2&=a_{1,2}^{-1}(Y_2+a_{1,3}Y_3+\cdots +a_{1,n}Y_n)&&\\
    Y_i&=X_i\ \forall\ i\in[n]\setminus\{2\}&&
  \end{align*}Then we get a quadratic form of the type $$Y_1Y_2+\sum_{2\leq i,j\leq n}c_{i,j}Y_iY_j$$ where $c_{i,j}\in\bbF_q$ for all $2\leq i,j\leq n$. So take the non-singular linear transformation   \begin{align*}
    Z_1& = Y_1+c_{2,2}Y_2+\cdots +c_{2,n}Y_n&&\\
    Z_i&=Y_i\ \forall\ i\in[n]\setminus\{1\}&&
  \end{align*}
  This transformation gives a quadratic form of the type $Z_1Z_2+g_3(Z_3,\dots, Z_n)$. So we have the lemma. 
\end{proof}

So now we will use this theorem to find proper canonical forms for any general quadratic forms. But we have to treat the cases of odd and even number of variables separately. 
\begin{Theorem}{Canonical Forms for even characteristic}{thm:qf-canonical-even-char}
  Let $f\in\bbF_q[X_1,\dots,X_n]$ be a non-degenerate quadratic form with $q$ even.
  \begin{enumerate}[label=(\roman*)]
    \item If $n$ is odd then $f$ is equivalent to $X_1X_2+X_3X_4+\cdots +X_{n-2}X_{n-1}+X_n^2$.
    \item If $n$ is even then $f$ is \begin{enumerate}
      \item either equivalent to $X_1X_2+\cdots+X_{n-1}X_n$, 
      \item or to $X_1X_2+\cdots+X_{n-3}X_{n-2}+X_{n-1}^2+a\cdot X_n^2$ where $a\in\bbF_q$ satisfies $\tr(a)=1$ where $\tr$ is the absolute trace function from $\bbF_q$ to $\bbF_2$. 
    \end{enumerate}
  \end{enumerate}
\end{Theorem}
\begin{proof}
  \begin{enumerate}[label=(\roman*)]
    \item Suppose $n$ is odd. By applying \QFRedEvenlem again and again on $f$ we get an equivalent quadratic form of the type $X_1X_2+X_3X_4+\cdots +X_{n-2}X_{n-1}+aX_n^2$ for some $a\in\bbF_q^*$. Then take the non-singular linear transformation
    \begin{align*}
      Y_n& = a^{-\nicefrac{q}{2}}X_n&&\\ 
      Y_i& = X_i\ \forall\ i\in[n-1]&&
    \end{align*}
    Then we get the desired quadratic form. 
    \item So now assume $n$ is even. Again by applying \QFRedEvenlem again and again on $f$ we get an equivalent quadratic form of the type $$X_1X_2+\cdots+X_{n-3}X_{n-2}+bX_{n-1}^2+cX_{n-1}X_n+dX_n^2$$with $b,c,d\in\bbF_q$. If $c=0$ then take the non-singular linear transformation
    \begin{align*}
      Y_{n-1}& = b^{\nicefrac{q}2}X_{n-1}+d^{\nicefrac{q}2}X_n&&\\ 
      Y_i& = X_i\ \forall\ i\in[n]\setminus\{n-1\}&&
    \end{align*}Then we get a quadratic form with $(n-1)$-variables. But $f$ is non-degenerate. Hence contradiction \ctr Therefore $c\in\bbF_q^*$. 
    
    If $b=0$ then take the non-singular linear transformation 
      \begin{align*}
      Y_{n-1}& = cX_{n-1}+dX_n&&\\ 
      Y_i& = X_i\ \forall\ i\in[n]\setminus\{n-1\}&&
    \end{align*}With this transformation we get the first form. 
    
    
    If $b\neq 0$ then take the non-singular linear transformation
    \begin{align*}
      Y_{n-1}& = b^{-\nicefrac{q}2}X_{n-1}&&\\
      Y_n & = b^{\nicefrac{q}2}c^{-1}X_n && \\ 
      Y_i& = X_i\ \forall\ i\in[n-2]&&
    \end{align*}Then we get a quadratic form of the type $$Y_1Y_2+\cdots+Y_{n-3}Y_{n-2}+Y_{n-1}^2+Y_{n-1}Y_n+aY_n^2$$for some $a\in\bbF_q$. If $X^2+X+a$ is reducible in $\bbF_q$ then there exists $c_1,c_2\in\bbF_q$ such that $$X^2+X+a=(X+c_1)(X+c_2)\implies Y_{n-1}^2+Y_{n-1}Y_n+aY_n^2=(Y_{n-1}+c_1Y_n)(Y_{n-1}+c_2Y_n)$$Therefore if we take the non-singular linear transformation 
    \begin{align*}
      Z_{n-1}& = Y_{n-1}+c_1Y_n&&\\
      Z_n & = Y_{n-1}+c_2Y_n && \\ 
      Z_i& = Y_i\ \forall\ i\in[n-2]&&
    \end{align*}then we get a quadratic form of the first form. Now $X^2+X+a$ is irreducible in $\bbF_q[X]$ if  $\tr(a)=1$. Then we have the second form. 
  \end{enumerate}
  Therefore we have the canonical forms for all cases. 
\end{proof}
Now by the above theorem it suffices to only find the number of solutions for canonical non-degenerate quadratic forms. We will first show two results on two variable version then we will use it to show the general version. 
\begin{lemma}{Solutions of a Two-Variable Quadratic Form in even characteristic}{lem:qf-two-var-count-even}
  For even $q$, let $a,\alpha\in\bbF_q$ with $\tr(a)=1$. Then $$Z\lt(X_1^2+X_1X_2+aX_2^2=\alpha\rt)=q-\vth(\alpha).$$where $\tr$ is the absolute trace function from $\bbF_q$ to $\bbF_2$. 
\end{lemma}
\begin{proof}
  Since $\tr(a)=1$ the polynomial $X^2+X+a$ is irrducible in $\bbF_q[X]$. So there exists $\gm\in\bbF_{q^2}\setminus\bbF_q$ such that $X^2+X+a=(X+\gm)(X+\gm^q)$. So $$f(X_1,X_2)=X_1^2+X_1X_2+aX_2^2=(X_1+\gm X_2)(X_1+\gm^q X_2)$$Therefore for any $(\beta_1,\beta_2)\in\bbF_q^2$ we have $$f(\beta_1,\beta_2)=(\beta_1+\gm\beta_2)(\beta_1+\gm^q\beta_2)=(\beta_1+\gm\beta_2)(\beta_1+\gm\beta_2)^q=(\beta_1+\gm\beta_2)^{q+1}$$So we get $$Z(f(X_1,X_2)=\alpha)=|\{\gm\in\bbF_{q^2}\colon \gm^{q+1}=\alpha\}|$$Hence if $\alpha=0$ then $Z\lt(X_1^2+X_1X_2+aX_2^2=0\rt)=1=q-\vth(0)$. If $\alpha\neq 0$ since $\bbF_{q^2}^*$ is cyclic we have $\alpha^{\nicefrac{(q^2-1)}{q+1}}=\alpha^{q-1}=1$ and hence there are $q+1$ elements in $\bbF_{q^2}$ such that $\gm^{q+1}=\alpha$. Hence $Z\lt(X_1^2+X_1X_2+aX_2^2=\alpha\rt)=q+1=q+\vth(\alpha)$.
\end{proof}
\begin{lemma}{}{lem:x1x2-quad-form}
  For even $q$, let $\alpha\in\bbF_q$. Then $Z(X_1X_2=\alpha)=q+\vth(\alpha)$. 
\end{lemma}
\begin{proof}
  If $\alpha\neq 0$ then both $X_1$ and $X_2$ has to be assigned with non-zero values. After assigning $X_1$ arbitrarily from $\bbF_q^*$ the value of $X_2$ gets uniquely determined. Therefore $Z(X_1X_2=\alpha)=q-1=q+\vth(\alpha)$. 

  Suppose $\alpha=0$. Then at least one of $X_1$ or $X_2$ gets assigned $0$. Then the other variable can take any value from $\bbF_q^*$. Therefore there are $2q-1$ such solutions. Hence $Z(X_1X_2=\alpha)=2q-1=q+\vth(\alpha)$. So we have the lemma. 
\end{proof}
\begin{Theorem}{Solution Counts for Canonical Quadratic Forms in even characteristic}{thm:qf-count-even-char}
  Let $\bbF_q$ be a finite field with $q$ even. Then for any $\alpha\in\bbF_q$:
  \begin{enumerate}[label=(\roman*)]
    \item If $n$ is odd, $Z\lt(X_1X_2+\cdots + X_{n-2}X_{n-1}+X_n^2=\alpha\rt)=q^{n-1}+\vth(\alpha)\cdot q^{\nicefrac{(n-1)}{2}}.$
    \item If $n$ is even:
    \begin{enumerate}[label=(\alph*)]
      \item $Z\lt(X_1X_2+\cdots +X_{n-1}X_n=\alpha\rt)=q^{n-1}+\vth(\alpha)\cdot q^{\nicefrac{(n-2)}{2}}$;
      \item for $a\in\bbF_q$ with $\tr(a)=1$, then $Z\lt(X_1X_2+\cdots+X_{n-1}X_{n}+X_{n-1}^2+a\cdot X_n^2=\alpha\rt)=q^{n-1}-\vth(\alpha)\cdot q^{\nicefrac{(n-2)}{2}}$, where $\tr$ is the absolute trace function from $\bbF_q$ to $\bbF_2$.
    \end{enumerate}
  \end{enumerate}
\end{Theorem}
\begin{proof}
  \begin{enumerate}[label=(\roman*)]
    \item Suppose $n$ is odd. Since $q$ is even the equation $X^2=\alpha$ for any $\alpha\in\bbF_q$ has an unique solution in $\bbF_q$. Therefore $Z\lt(X_1X_2+\cdots + X_{n-2}X_{n-1}+X_n^2=\alpha\rt)=q^{n-1}+\vth(\alpha)\cdot q^{\nicefrac{(n-1)}{2}}$ because after assigning the variables $X_1,\dots, X_{n-1}$ arbitrarily the value of $X_n$ is uniquely determined. 
    \item So let $n$ is even. Let $m=\nicefrac{n}2$. Then we have \begin{align*}
        Z\lt(X_1X_2+\cdots +X_{n-1}X_n=\alpha\rt)& = \sum_{\substack{\alpha_1,\dots,\alpha_m\in\bbF_q\\ \alpha_1+\cdots+\alpha_m=\alpha}}\prod_{j=1}^n Z(X_{2j-1}X_{2j}=\alpha_j)\\ 
        & = \sum_{\substack{\alpha_1,\dots,\alpha_m\in\bbF_q\\ \alpha_1+\cdots+\alpha_m=\alpha}} \prod_{j=1}^m (q+\vth(\alpha_j)) \byl{x1x2-quad-form}\\ 
        & = q^{2m-1}+\sum_{\substack{\alpha_1,\dots,\alpha_m\in\bbF_q\\ \alpha_1+\cdots+\alpha_m=\alpha}}\prod_{j=1}^m\vth(\alpha_j)\\ 
        & = q^{2m-1}+\vth(\alpha)\cdot q^{m-1}=q^{n-1}+\vth(\alpha)\cdot q^{\nicefrac{(n-2)}{2}}\by{\QFConvProplem}
      \end{align*}
      So we have the first part. Now for the second part we already showed the $n=2$ case in \QFTwoVarEvenlem. So assume $n\geq 4$.  
      \begin{align*}
        & Z\lt(X_1X_2+\cdots+X_{n-3}X_{n-2}+X_{n-1}^2+a\cdot X_n^2=\alpha\rt)\\ 
        = \;&\;\sum_{\substack{\alpha_1,\alpha_2\in\bbF_q\\ \alpha_1+\alpha_2=\alpha}}Z\lt(X_1X_2+\cdots+X_{n-3}X_{n-2}=\alpha_1\rt)\cdot Z(X_{n-1}X_{n}+X_{n-1}^2+a\cdot X_n^2=\alpha_2)\\ 
        = \;&\; \sum_{\substack{\alpha_1,\alpha_2\in\bbF_q\\ \alpha_1+\alpha_2=\alpha}} \lt(q^{n-3}+\vth(\alpha_1)\cdot q^{\nicefrac{(n-4)}{2}}\rt)(q-\vth(\alpha_2))\by{first part and \QFTwoVarEvenlem}\\ 
        = \;&\; q^{n-1}-q^{\nicefrac{(n-4)}{2}}\sum_{\substack{\alpha_1,\alpha_2\in\bbF_q\\ \alpha_1+\alpha_2=\alpha}}\vth(\alpha_1)\cdot\vth(\alpha_2) \by{\QFConvProplem}\\ 
        = \;&\; q^{n-1}-q^{\nicefrac{(n-4)}{2}}\cdot q\cdot \vth(\alpha)\\ 
        = \;&\; q^{n-1}-\vth(\alpha)\cdot q^{\nicefrac{(n-2)}{2}}
      \end{align*}
      So now we have the result for second part
  \end{enumerate}
  Therefore we get the values of number of solutions of quadratic forms in even characteristic finite fields. 
\end{proof}